\documentclass[11pt,a4paper]{article}

\usepackage[T1]{fontenc}
\usepackage[utf8]{inputenc}
\usepackage{lmodern}
\usepackage[margin=1in]{geometry}
\usepackage{amsmath,amssymb,amsthm}
\usepackage{booktabs}
\usepackage{listings}
\usepackage{microtype}
\usepackage{xcolor}
\usepackage[hidelinks]{hyperref}

\hypersetup{
  pdftitle={Formal Verification of a Collatz Residue Sieve},
  pdfsubject={Number theory, formal verification},
  colorlinks=true, linkcolor=black, citecolor=black, urlcolor=blue!55!black
}

% ---------------------------------------------------------------- theorems --
\theoremstyle{plain}
\newtheorem{theorem}{Theorem}[section]
\newtheorem{lemma}[theorem]{Lemma}
\newtheorem{proposition}[theorem]{Proposition}
\newtheorem{corollary}[theorem]{Corollary}
\theoremstyle{definition}
\newtheorem{definition}[theorem]{Definition}
\newtheorem{remark}[theorem]{Remark}

% ------------------------------------------------------------- pseudocode ---
\newcounter{algline}
\newenvironment{pseudocode}
  {\begin{list}{\footnotesize\sffamily\arabic{algline}:}%
     {\usecounter{algline}%
      \setlength{\leftmargin}{2.4em}\setlength{\labelwidth}{1.6em}%
      \setlength{\labelsep}{0.8em}\setlength{\itemsep}{1pt}%
      \setlength{\parsep}{0pt}\setlength{\topsep}{4pt}}}
  {\end{list}}
\newcommand{\kw}[1]{\textbf{#1}}
\newcommand{\ind}{\hspace*{1.4em}}
\newcommand{\indd}{\hspace*{2.8em}}
\newcommand{\comm}[1]{{\color{gray}\(\triangleright\) \textit{#1}}}

% -------------------------------------------------------------- listings ----
\lstdefinestyle{py}{
  basicstyle=\footnotesize\ttfamily, keywordstyle=\color{blue!60!black},
  commentstyle=\color{gray}\itshape, stringstyle=\color{green!40!black},
  numbers=left, numberstyle=\tiny\color{gray}, numbersep=8pt,
  frame=leftline, framesep=6pt, rulecolor=\color{gray!40},
  showstringspaces=false, breaklines=true, language=Python, xleftmargin=1.2em
}
\lstdefinelanguage{Rocq}{
  morekeywords={Definition,Fixpoint,Theorem,Lemma,Corollary,Proof,Qed,forall,
    exists,match,with,end,if,then,else,let,in,fun,Require,Import,induction,
    intros,Example,Print,Assumptions,Local,Open,Scope},
  morecomment=[n]{(*}{*)}, sensitive=true
}
\lstdefinestyle{rocq}{
  basicstyle=\footnotesize\ttfamily, keywordstyle=\color{purple!70!black},
  commentstyle=\color{gray}\itshape, frame=leftline, framesep=6pt,
  rulecolor=\color{gray!40}, breaklines=true, language=Rocq, xleftmargin=1.2em
}

% ---------------------------------------------------------------- macros ----
\newcommand{\N}{\mathbb{N}}
\newcommand{\Np}{\mathbb{N}_{\geq 1}}
\newcommand{\Tmap}{\mathrm{T}}
\newcommand{\vtwo}{v_2}
\newcommand{\Acoef}{A}
\newcommand{\Ccoef}{C}
\newcommand{\reaches}{\mathsf{reaches1}}
\newcommand{\Rlev}{\mathcal{R}}
\newcommand{\Slev}{\mathcal{S}}

\begin{document}

% =============================================================== title ======
\title{\bf Formal Verification of a Collatz Residue Sieve\\[2mm]
\large A machine-checked correctness proof of \texttt{Residual\_Analysis.py}}

\author{RobinCodes\\[1mm]
\small\href{https://github.com/RobinCodes/Collatz}{\texttt{github.com/RobinCodes/Collatz}}}

\date{17 August 2026}

\maketitle

\begin{abstract}
\noindent
We give a complete correctness proof, machine-checked in the Rocq Prover, of
\texttt{Residual\_Analysis.py} --- a program that iteratively sieves residue
classes modulo $2^q$ that cannot contain a minimal counterexample to the
Collatz conjecture. Three aspects of the implementation are non-obvious and
are the substance of the proof. First, the inner loop truncates its final
right-shift at the step budget, producing an intermediate value that is not an
element of the Collatz trajectory; we show it nevertheless equals
$\Tmap^{q}(n)$ exactly, where $\Tmap$ is the accelerated Collatz map. Second,
the program tests only the least positive representative of each residue
class; we prove via a constructive form of Terras' theorem that this single
test decides the whole class. Third, the program treats a fixed point
$\Tmap^{q}(r)=r$ as a \emph{dead} class while separately recording it; we show
this is sound, because every strictly larger member of such a class descends.
The main theorem states that for every odd $n \geq 1$ whose class the sieve has
discarded, either $n$ reaches $1$ given that all smaller positive integers do
--- so $n$ is not a minimal counterexample --- or $n$ lies on a cycle, and
every such $n$ is reported in the program's \texttt{cycles} list. The Rocq
development is axiom-free and contains no admitted proofs.
\end{abstract}

% =============================================================== intro ======
\section{Introduction}

The Collatz conjecture asserts that iterating
\[
  n \;\longmapsto\;
  \begin{cases}
    n/2 & n \text{ even},\\
    3n+1 & n \text{ odd},
  \end{cases}
\]
from any positive integer eventually reaches $1$. It remains open; computational
searches have confirmed it for all $n$ below roughly $2^{68}$.

A standard partial attack is the \emph{residue sieve} of Terras~\cite{terras}
and Everett~\cite{everett}. Because the first $q$ steps of the accelerated map
depend only on $n \bmod 2^{q}$, one may partition $\Np$ into $2^{q}$ residue
classes and ask, for each class, whether \emph{every} member drops below
itself within $q$ steps. A class with that property can contain no minimal
counterexample and may be discarded; a surviving class is refined modulo
$2^{q+1}$ and retested. Terras proved that the density of surviving classes
tends to $0$; the conjecture nevertheless remains open, because a density-zero
set may still be nonempty.

The program \texttt{Residual\_Analysis.py} in the repository
\href{https://github.com/RobinCodes/Collatz}{\texttt{github.com/RobinCodes/Collatz}}
implements this sieve. Its arithmetic core is compact and, in three places,
does something that looks wrong at first reading. This paper proves that it is
right, and states precisely what it does and does not establish.

\paragraph{Contributions.}
\begin{enumerate}\itemsep2pt
\item A precise specification of the algorithm (Section~\ref{sec:prog}) and of
      the accelerated map's affine structure (Section~\ref{sec:affine}).
\item A constructive \emph{shift lemma} (Theorem~\ref{thm:shift}) replacing the
      usual congruence formulation of Terras' theorem, which is what makes the
      Rocq development tractable.
\item Proofs that the truncated shift loop computes $\Tmap^{q}$
      (Proposition~\ref{prop:eval}), that the least-representative test decides
      the class (Theorem~\ref{thm:descent}), and that the sieve discards no
      potential minimal counterexample (Theorem~\ref{thm:main}).
\item A full Rocq formalization (Section~\ref{sec:rocq}), axiom-free, with the
      structural invariants that justify the program's reported density figure.
\item An explicit account of the algorithm's limits (Section~\ref{sec:limits}).
\end{enumerate}

% ============================================================ definitions ===
\section{The accelerated map}
\label{sec:defs}

\begin{definition}[Accelerated Collatz map]
$\Tmap : \N \to \N$ is
\[
  \Tmap(n) \;=\;
  \begin{cases}
    n/2 & n \text{ even},\\[1mm]
    (3n+1)/2 & n \text{ odd}.
  \end{cases}
\]
We write $\Tmap^{q}$ for the $q$-fold iterate and set
$\reaches(n) \iff \exists k.\, \Tmap^{k}(n) = 1$.
\end{definition}

Since $3n+1$ is even for odd $n$, both branches are exact integer divisions.
The following triviality is the reason the program can count \emph{halvings}
rather than \emph{steps}.

\begin{lemma}[Step accounting]\label{lem:accounting}
Every application of $\Tmap$ divides exactly once by $2$. Consequently, for
any $n$ and $q$, the value obtained from $n$ by repeatedly applying
$m \mapsto 3m+1$ at odd values and dividing out powers of two, stopping the
instant the \emph{total} number of divisions by $2$ reaches $q$, is
$\Tmap^{q}(n)$.
\end{lemma}

\begin{definition}[$2$-adic valuation]
For $y \neq 0$, $\vtwo(y)$ is the largest $e$ with $2^{e} \mid y$.
\end{definition}

\begin{lemma}[Chained halving]\label{lem:chain}
Let $m$ be odd and let $d \geq 1$ satisfy $2^{d} \mid 3m+1$. Then
\[
  \Tmap^{d}(m) \;=\; \frac{3m+1}{2^{d}} .
\]
\end{lemma}

\begin{proof}
Induction on $d$. For $d=1$: $m$ is odd, so $\Tmap(m) = (3m+1)/2$. Suppose the
claim for $d$ and let $2^{d+1} \mid 3m+1$; then also $2^{d} \mid 3m+1$, so
$x := \Tmap^{d}(m) = (3m+1)/2^{d}$ by the induction hypothesis. Writing
$3m+1 = 2^{d+1}t$ gives $x = 2t$, which is even, hence
$\Tmap^{d+1}(m) = \Tmap(x) = x/2 = (3m+1)/2^{d+1}$.
\end{proof}

Note that Lemma~\ref{lem:chain} does \emph{not} require $d = \vtwo(3m+1)$. This
is exactly the freedom the implementation exploits.

% =============================================================== program ====
\section{The program}
\label{sec:prog}

\subsection{Source}

The arithmetic core of \texttt{Residual\_Analysis.py} is:

\begin{lstlisting}[style=py]
def collatz_step(n: int, score: int) -> tuple[int, int]:
	y = 3 * n + 1
	v2 = (y & -y).bit_length() - 1

	d = min(v2, q - score)
	return y >> d, d

def collatz_sim(n: int) -> bool:
	og_n = n
	score = 0
	while score < q:
		data = collatz_step(n, score)
		n = data[0]
		score += data[1]

	if og_n < n:
		return True   # Survives
	elif og_n > n:
		return False  # Dies
	else:
		cycles.append(og_n)
		return False
\end{lstlisting}

\subsection{Pseudocode}

We abstract the two loops. Algorithm~\ref{alg:eval} is
\texttt{collatz\_sim}'s inner loop; Algorithm~\ref{alg:sieve} is the outer
sieve, including the reporting behaviour of the current version of the
program.

\begin{figure}[htb]
\hrule\vspace{4pt}
\noindent\textbf{Algorithm 1}\; $\textsc{Eval}(n, q)$ ---
evaluate $\Tmap^{q}(n)$ for odd $n$ by truncated shifting.
\vspace{2pt}\hrule
\begin{pseudocode}
\item $\mathit{score} \gets 0$; \quad $m \gets n$
\item \kw{while} $\mathit{score} < q$ \kw{do}
\item \ind $y \gets 3m + 1$ \hfill \comm{$y$ is even, since $m$ is odd}
\item \ind $d \gets \min\!\left(\vtwo(y),\; q - \mathit{score}\right)$
        \hfill \comm{$1 \leq d \leq q - \mathit{score}$}
\item \ind $m \gets y \,/\, 2^{d}$ \hfill \comm{\texttt{y >> d}}
\item \ind $\mathit{score} \gets \mathit{score} + d$
\item \kw{return} $m$
\end{pseudocode}
\hrule
\caption{}
\label{alg:eval}
\end{figure}

\begin{figure}[htb]
\hrule\vspace{4pt}
\noindent\textbf{Algorithm 2}\; $\textsc{Sieve}$ --- iterative residue
elimination.
\vspace{2pt}\hrule
\begin{pseudocode}
\item $q \gets 1$; \quad $\Rlev \gets [\,1\,]$; \quad $\mathcal{C} \gets [\;]$
\item \kw{loop}
\item \ind $\Slev \gets [\;]$
\item \ind \kw{for} $r \in \Rlev$ \kw{do}
\item \indd $m \gets \textsc{Eval}(r, q)$
\item \indd \kw{if} $m > r$ \kw{then} append $r$ to $\Slev$
        \hfill \comm{class survives}
\item \indd \kw{else if} $m = r$ \kw{then} append $r$ to $\mathcal{C}$
        \hfill \comm{cycle element; class still dies}
\item \ind \kw{report} $\;100 - 100\,|\Slev| / 2^{q}$, \; $q$, \;
        $\Slev$ if $q \leq 15$, \; $\mathcal{C}$
\item \ind $\Rlev \gets \bigl[\, r,\; r + 2^{q} \;:\; r \in \Slev \,\bigr]$
\item \ind $q \gets q + 1$
\end{pseudocode}
\hrule
\caption{}
\label{alg:sieve}
\end{figure}

Three features deserve comment, and each is discharged below.

\begin{enumerate}\itemsep3pt
\item[(P1)] Line 3 of Algorithm~\ref{alg:eval} applies $3m+1$
      unconditionally, with no parity test. This is safe only if $m$ is odd on
      every entry to the loop body, which is not immediate: line~5 can leave
      $m$ \emph{even} when the budget truncates the shift.
\item[(P2)] Line 4 caps $d$ at the remaining budget. When
      $d < \vtwo(y)$ the returned $m$ is not a Collatz trajectory element at
      all. The comparison on line~6 of Algorithm~\ref{alg:sieve} is
      nevertheless meaningful.
\item[(P3)] Only the least representative $r$ of each class is ever tested,
      yet the conclusion drawn concerns the entire infinite class
      $\{\, r + j 2^{q} : j \geq 0 \,\}$.
\end{enumerate}

\begin{proposition}[Loop invariant and (P1)]\label{prop:inv}
Let $n$ be odd and $q \geq 0$. Throughout Algorithm~\ref{alg:eval},
\[
  m = \Tmap^{\mathit{score}}(n)
  \qquad\text{and}\qquad
  \bigl(\, m \text{ is odd} \;\lor\; \mathit{score} = q \,\bigr).
\]
Moreover $d \geq 1$ at every iteration, so the loop terminates after at most
$q$ iterations.
\end{proposition}

\begin{proof}
Initially $m = n = \Tmap^{0}(n)$ is odd. Assume the invariant with
$\mathit{score} < q$; then $m$ is odd, so $y = 3m+1$ is even and
$\vtwo(y) \geq 1$, whence $d = \min(\vtwo(y), q - \mathit{score}) \geq 1$ and
$d \leq \vtwo(y)$, so $2^{d} \mid y$. By Lemma~\ref{lem:chain},
$y/2^{d} = \Tmap^{d}(m) = \Tmap^{\mathit{score}+d}(n)$, establishing the first
conjunct. For the second: if $d = \vtwo(y)$ then $y/2^{d}$ is odd by
maximality of the valuation; otherwise $d = q - \mathit{score}$, and the new
score equals $q$. Since $d \geq 1$, the score strictly increases.
\end{proof}

\begin{proposition}[(P2)]\label{prop:eval}
For odd $n$ and every $q \geq 0$, $\textsc{Eval}(n,q) = \Tmap^{q}(n)$.
\end{proposition}

\begin{proof}
Immediate from Proposition~\ref{prop:inv}: the loop exits with
$\mathit{score} = q$ --- the score increases by $d \leq q - \mathit{score}$, so
it never overshoots --- and the invariant then reads $m = \Tmap^{q}(n)$.
\end{proof}

Proposition~\ref{prop:eval} is the resolution of the truncation puzzle. The
final value may be even and may not occur anywhere in the Collatz trajectory
of $n$; it is nonetheless \emph{exactly} $\Tmap^{q}(n)$, because the
truncated division is the tail of the $q$-th $\Tmap$-step, not a corruption of
it.

% ================================================================ affine ====
\section{Affine structure}
\label{sec:affine}

\begin{definition}[Affine coefficients]\label{def:ac}
Define $\Acoef_{q}(n), \Ccoef_{q}(n) \in \N$ by recursion on $q$:
\[
  (\Acoef_{0}, \Ccoef_{0}) = (1, 0), \qquad
  (\Acoef_{q+1}, \Ccoef_{q+1}) =
  \begin{cases}
    (\Acoef_{q},\; \Ccoef_{q})
      & \Tmap^{q}(n) \text{ even},\\[1mm]
    (3\Acoef_{q},\; 3\Ccoef_{q} + 2^{q})
      & \Tmap^{q}(n) \text{ odd}.
  \end{cases}
\]
\end{definition}

\begin{lemma}[Affine identity]\label{lem:affine}
For all $n, q$:
\[
  2^{q}\,\Tmap^{q}(n) \;=\; \Acoef_{q}(n)\, n \;+\; \Ccoef_{q}(n),
  \qquad \Acoef_{q}(n) = 3^{a_{q}(n)} \text{ is odd},
  \qquad \Ccoef_{q}(n) \geq 0,
\]
where $a_{q}(n)$ is the number of odd steps among the first $q$.
\end{lemma}

\begin{proof}
Induction on $q$. The base case is $n = n$. Write $m = \Tmap^{q}(n)$ and assume
$2^{q}m = \Acoef_{q}n + \Ccoef_{q}$.
If $m$ is even, $\Tmap^{q+1}(n) = m/2$ and
$2^{q+1}(m/2) = 2^{q}m = \Acoef_{q}n + \Ccoef_{q}$, matching the unchanged
coefficients.
If $m$ is odd, $\Tmap^{q+1}(n) = (3m+1)/2$ and
$2^{q+1}\cdot\frac{3m+1}{2} = 2^{q}(3m+1) = 3(\Acoef_{q}n + \Ccoef_{q}) + 2^{q}
= (3\Acoef_{q})n + (3\Ccoef_{q} + 2^{q})$.
Oddness of $\Acoef_{q}$ is immediate: $1$ is odd and $3\Acoef_{q}$ is odd
whenever $\Acoef_{q}$ is.
\end{proof}

The next theorem is Terras' theorem in the form we need. The usual statement
is a congruence --- if $n \equiv n' \pmod{2^{q}}$ then the first $q$ parities
agree --- but the constructive version below is both stronger and easier to
formalize, since it computes the displacement rather than merely asserting
equality of parity vectors.

\begin{theorem}[Shift lemma]\label{thm:shift}
For all $q \in \N$ and all $n, j \in \N$,
\[
  \Tmap^{q}\!\left(n + j\,2^{q}\right) \;=\; \Tmap^{q}(n) + \Acoef_{q}(n)\, j,
  \qquad\text{and}\qquad
  \bigl(\Acoef_{q},\Ccoef_{q}\bigr)\!\left(n + j\,2^{q}\right)
    = \bigl(\Acoef_{q},\Ccoef_{q}\bigr)(n).
\]
\end{theorem}

\begin{proof}
Simultaneous induction on $q$. For $q=0$ both claims are trivial
($\Acoef_{0} = 1$).

For the step, put $P = 2^{q}$ and note $n + j\,2^{q+1} = n + (2j)P$. Write
$A = \Acoef_{q}(n)$ and $m = \Tmap^{q}(n)$. The induction hypothesis applied
with displacement $2j$ gives
\[
  \Tmap^{q}\!\left(n + j\,2^{q+1}\right) = m + 2Aj,
  \qquad
  \bigl(\Acoef_{q},\Ccoef_{q}\bigr)\!\left(n + j\,2^{q+1}\right)
    = \bigl(\Acoef_{q},\Ccoef_{q}\bigr)(n).
\]
Since $2Aj$ is even, $m + 2Aj$ has the same parity as $m$; hence the
coefficient recursion of Definition~\ref{def:ac} takes the same branch at
step $q+1$ for both arguments, which proves the second claim. For the first:
\begin{itemize}\itemsep2pt
\item $m$ even: $\Tmap(m + 2Aj) = \dfrac{m + (Aj)\cdot 2}{2}
      = \dfrac{m}{2} + Aj = \Tmap^{q+1}(n) + \Acoef_{q+1}(n)\,j$,
      using $\Acoef_{q+1}(n) = A$.
\item $m$ odd: $m + 2Aj$ is odd and
      $\Tmap(m + 2Aj) = \dfrac{3m + 1 + (3Aj)\cdot 2}{2}
      = \dfrac{3m+1}{2} + 3Aj = \Tmap^{q+1}(n) + \Acoef_{q+1}(n)\,j$,
      using $\Acoef_{q+1}(n) = 3A$.
\end{itemize}
In both cases the identity $\left(x + y\cdot 2\right)/2 = x/2 + y$ for exact
integer division is all that is required.
\end{proof}

\begin{corollary}[Class uniformity]\label{cor:uniform}
$\Tmap^{q}$ restricted to a residue class modulo $2^{q}$ is the affine map
$r + j2^{q} \mapsto \Tmap^{q}(r) + \Acoef_{q}(r)\,j$. In particular
$\Tmap^{q}(n) - n$ is monotone in $j$: strictly decreasing if
$\Acoef_{q}(r) < 2^{q}$ and strictly increasing if $\Acoef_{q}(r) > 2^{q}$.
\end{corollary}

% =============================================================== descent ====
\section{The least-representative test}
\label{sec:descent}

We now discharge (P3): why testing $r$ alone suffices.

\begin{lemma}\label{lem:Asmall}
Let $q \geq 1$ and $r \geq 1$ with $\Tmap^{q}(r) \leq r$. Then
$\Acoef_{q}(r) < 2^{q}$.
\end{lemma}

\begin{proof}
By Lemma~\ref{lem:affine},
$\Acoef_{q}(r)\,r + \Ccoef_{q}(r) = 2^{q}\Tmap^{q}(r) \leq 2^{q}r$.
Since $\Ccoef_{q}(r) \geq 0$ we get $\Acoef_{q}(r)\,r \leq 2^{q}r$, and since
$r \geq 1$, $\Acoef_{q}(r) \leq 2^{q}$. Equality is impossible:
$\Acoef_{q}(r)$ is odd by Lemma~\ref{lem:affine}, while $2^{q}$ is even for
$q \geq 1$. Hence $\Acoef_{q}(r) < 2^{q}$.
\end{proof}

\begin{theorem}[Descent]\label{thm:descent}
Let $q \geq 1$ and $r \geq 1$ with $\Tmap^{q}(r) \leq r$. Then for every
$j \geq 0$ such that $\Tmap^{q}(r) < r$ or $j \geq 1$,
\[
  \Tmap^{q}\!\left(r + j\,2^{q}\right) \;<\; r + j\,2^{q}.
\]
\end{theorem}

\begin{proof}
Write $A = \Acoef_{q}(r)$; by Lemma~\ref{lem:Asmall}, $A < 2^{q}$. By
Theorem~\ref{thm:shift}, $\Tmap^{q}(r + j2^{q}) = \Tmap^{q}(r) + Aj$.
If $\Tmap^{q}(r) < r$ then $\Tmap^{q}(r) + Aj < r + 2^{q}j$ since
$Aj \leq 2^{q}j$. If instead $j \geq 1$ then $Aj < 2^{q}j$ strictly, and
$\Tmap^{q}(r) \leq r$ suffices.
\end{proof}

Theorem~\ref{thm:descent} covers both branches the program takes:

\begin{itemize}\itemsep3pt
\item \textbf{Strict descent} ($\Tmap^{q}(r) < r$): the whole class descends,
      including $r$ itself. The program discards it.
\item \textbf{Fixed point} ($\Tmap^{q}(r) = r$): every member \emph{strictly
      above} $r$ descends, but $r$ itself does not --- it lies on a cycle. The
      program discards the class and separately records $r$ in
      \texttt{cycles}. This is sound, and it is why the equality branch is not
      a bug.
\end{itemize}

\begin{remark}
The fixed-point case has a pleasant reformulation. If $\Tmap^{q}(r) = r$ then
Lemma~\ref{lem:affine} gives $\Ccoef_{q}(r) = (2^{q} - \Acoef_{q}(r))\,r$, so
for $n = r + j2^{q}$,
\[
  \Tmap^{q}(n) - n \;=\; \frac{(2^{q} - \Acoef_{q}(r))\,(r - n)}{2^{q}} \;<\; 0
  \qquad (j \geq 1),
\]
an explicit descent rate.
\end{remark}

\begin{theorem}[No minimal counterexample]\label{thm:nomin}
Let $q \geq 1$, $r \geq 1$, $\Tmap^{q}(r) \leq r$, and let
$n = r + j2^{q}$ with $\Tmap^{q}(r) < r$ or $j \geq 1$. If every $m$ with
$1 \leq m < n$ satisfies $\reaches(m)$, then $\reaches(n)$.
\end{theorem}

\begin{proof}
By Theorem~\ref{thm:descent}, $\Tmap^{q}(n) < n$. Since $n \geq 1$ and
$\Tmap$ maps nonzero values to nonzero values, $\Tmap^{q}(n) \geq 1$. By
hypothesis $\reaches(\Tmap^{q}(n))$, say $\Tmap^{k}(\Tmap^{q}(n)) = 1$; then
$\Tmap^{k+q}(n) = 1$, so $\reaches(n)$.
\end{proof}

% ================================================================= sieve ====
\section{The sieve}
\label{sec:sieve}

\begin{definition}\label{def:sieve}
Let $\mathrm{survives}(q, r) \iff \Tmap^{q}(r) > r$ and, for a list $L$,
\[
  \mathrm{children}_{q}(L) \;=\;
  \bigl[\; r,\; r + 2^{q} \;:\; r \in L \;\bigr].
\]
Define $\Rlev_{1} = [\,1\,]$ and
$\Rlev_{q+1} = \mathrm{children}_{q}\bigl(\mathrm{filter}\,
\mathrm{survives}(q,\cdot)\;\Rlev_{q}\bigr)$.
\end{definition}

$\Rlev_{q}$ is exactly the program's \texttt{relevant\_residues} at level $q$,
and $\Slev_{q} = \mathrm{filter}\,\mathrm{survives}(q,\cdot)\,\Rlev_{q}$ is
\texttt{surviving\_residues}.

\begin{proposition}[Structural invariants]\label{prop:struct}
For every $q \geq 1$ and every $r \in \Rlev_{q}$:
\begin{enumerate}\itemsep1pt
\item[(i)] $r$ is odd;
\item[(ii)] $1 \leq r < 2^{q}$, so $r$ is the least positive member of its
      class modulo $2^{q}$;
\item[(iii)] $\Rlev_{q}$ has no repeated entries.
\end{enumerate}
\end{proposition}

\begin{proof}
Induction on $q$. For $q=1$, $\Rlev_{1} = [1]$ and $1 < 2$.
(i) $2^{q}$ is even for $q \geq 1$, so $r$ and $r + 2^{q}$ have the same
parity. (ii) If $r < 2^{q}$ then both $r$ and $r + 2^{q}$ lie below
$2^{q+1}$. (iii) Suppose $\Rlev_{q}$ is repetition-free with all entries below
$2^{q}$. Within one parent $r$ we have $r \neq r + 2^{q}$; across distinct
parents $r \neq r'$, we have $r \neq r'$, $r + 2^{q} \neq r' + 2^{q}$, and
$r \neq r' + 2^{q}$ because $r < 2^{q} \leq r' + 2^{q}$.
\end{proof}

Part (i) justifies (P1) at the level of the outer loop: every value fed to
\textsc{Eval} is odd, so Proposition~\ref{prop:inv} applies. Part (ii)
justifies (P3): the tested value really is the least member of its class, and
by Corollary~\ref{cor:uniform} that is the hardest case. Part (iii) is what
makes the program's \texttt{len(set(surviving\_residues))} equal to
\texttt{len(surviving\_residues)} --- the \texttt{set()} is redundant but
harmless --- and it is what makes the reported density well defined.

\begin{corollary}[The reported density]\label{cor:density}
The quantity $100 - 100\,|\Slev_{q}|/2^{q}$ printed at level $q$ is exactly the
percentage of residue classes modulo $2^{q}$ that the sieve has proven to
descend. All even classes are eliminated at $q=1$, since
$\Tmap(2m) = m < 2m$, which is why the denominator is $2^{q}$ rather than the
count $2^{q-1}$ of odd classes.
\end{corollary}

\begin{remark}
Corollary~\ref{cor:density} settles which of two natural quantities the
program reports. Relative to the odd residues alone --- the only ones the
sieve ever refines --- the surviving density is twice
$|\Slev_{q}|/2^{q}$. Both are meaningful; the printed figure is the density
among \emph{all} integers.
\end{remark}

\begin{theorem}[Coverage]\label{thm:cov}
Let $n \geq 1$ be odd and let $q \geq 1$. Then at least one of the following
holds:
\begin{enumerate}\itemsep1pt
\item[(a)] there are $r \in \Rlev_{q}$ and $j \geq 0$ with $n = r + j2^{q}$; or
\item[(b)] there are $q' \leq q$, $r \in \Rlev_{q'}$ and $j \geq 0$ with
      $n = r + j 2^{q'}$ and $\Tmap^{q'}(r) \leq r$.
\end{enumerate}
\end{theorem}

\begin{proof}
Induction on $q$. For $q=1$: $n$ is odd, so $n = 1 + j\cdot 2$ for some $j$,
and $1 \in \Rlev_{1}$, giving (a).

Assume the claim at $q$. If (b) holds, it holds verbatim at $q+1$. Otherwise
take $r \in \Rlev_{q}$ and $j$ with $n = r + j2^{q}$.
If $\mathrm{survives}(q,r)$ fails, then $\Tmap^{q}(r) \leq r$ and we are in
case (b) with $q' = q$. If it holds, then $r \in \Slev_{q}$, so both $r$ and
$r + 2^{q}$ belong to $\Rlev_{q+1}$. Split on the parity of $j$:
\[
  j = 2j' \;\Rightarrow\; n = r + j'2^{q+1},
  \qquad
  j = 2j' + 1 \;\Rightarrow\; n = (r + 2^{q}) + j'2^{q+1},
\]
which is case (a) at level $q+1$.
\end{proof}

\begin{theorem}[Main theorem: soundness of the sieve]\label{thm:main}
Let $n \geq 1$ be odd, let $q \geq 1$, and suppose every $m$ with
$1 \leq m < n$ satisfies $\reaches(m)$. Then at least one of the following
holds:
\begin{enumerate}\itemsep1pt
\item[(a)] $n$'s residue class is still present in $\Rlev_{q}$ --- the sieve
      makes no claim about $n$; or
\item[(b)] $\reaches(n)$ --- so $n$ is not a minimal counterexample; or
\item[(c)] $\Tmap^{q'}(n) = n$ for some $q' \leq q$, i.e.\ $n$ lies on a
      cycle of the accelerated map.
\end{enumerate}
Moreover, in case (c) the value $n$ appears in the program's \texttt{cycles}
list.
\end{theorem}

\begin{proof}
Apply Theorem~\ref{thm:cov}. Case (a) is immediate. In case (b) of that
theorem, we have $n = r + j2^{q'}$ with $r \geq 1$ and
$\Tmap^{q'}(r) \leq r$.

If $j \geq 1$, Theorem~\ref{thm:nomin} yields $\reaches(n)$.
If $j = 0$ then $n = r$. Either $\Tmap^{q'}(n) < n$, in which case
$\Tmap^{q'}(n)$ is a positive integer below $n$, so $\reaches(\Tmap^{q'}(n))$
by hypothesis and hence $\reaches(n)$; or $\Tmap^{q'}(n) = n$, which is case
(c).

For the final claim, observe that the coverage proof produces case (b) only
with $r \in \Rlev_{q'}$, and the program tests every element of $\Rlev_{q'}$
at level $q'$ and appends precisely those with $\Tmap^{q'}(r) = r$ to
\texttt{cycles}. Since $q' \leq q$, the value has already been recorded by the
time level $q$ is reported.
\end{proof}

\begin{corollary}
If the sieve at level $q$ discards a residue class, that class contains no
minimal counterexample to the Collatz conjecture other than, possibly, a cycle
element that the program reports.
\end{corollary}

% ================================================================== rocq ====
\section{Rocq formalization}
\label{sec:rocq}

The development \texttt{proofs/CollatzSieve.v} formalizes everything above.
Numbers are binary naturals (\texttt{N}); step counts are \texttt{nat}, so
that iteration is structurally recursive.

\begin{lstlisting}[style=rocq]
Definition T (n : N) : N := if N.even n then n / 2 else (3 * n + 1) / 2.

Fixpoint Tp (k : nat) (n : N) : N :=
  match k with O => n | S k' => T (Tp k' n) end.

Definition reaches1 (n : N) : Prop := exists k, Tp k n = 1.
\end{lstlisting}

The inner loop is modelled with an explicit structural fuel argument, since
$\mathit{rem} - d$ is not syntactically decreasing:

\begin{lstlisting}[style=rocq]
Fixpoint sim_aux (fuel rem : nat) (n : N) : N :=
  match fuel with
  | O   => n
  | S f => match rem with
           | O   => n
           | S _ => let y := 3 * n + 1 in
                    let d := Nat.min (v2 y) rem in
                    sim_aux f (rem - d) (y / pow2 d)
           end
  end.

Definition sim (q : nat) (n : N) : N := sim_aux q q n.
\end{lstlisting}

Since $d \geq 1$ at every iteration, fuel $= q$ always suffices; this is the
content of the hypothesis $\mathit{rem} \leq \mathit{fuel}$ in
\texttt{sim\_aux\_correct}.

The $2$-adic valuation is defined by structural recursion on the binary
representation, which makes it total and evaluation-friendly:

\begin{lstlisting}[style=rocq]
Fixpoint v2p (p : positive) : nat :=
  match p with xO q => S (v2p q) | _ => O end.

Definition v2 (y : N) : nat := match y with N0 => O | Npos p => v2p p end.
\end{lstlisting}

Table~\ref{tab:map} maps the results of this paper to their Rocq names.

\begin{table}[htb]
\centering
\small
\begin{tabular}{@{}lll@{}}
\toprule
\textbf{Result} & \textbf{Statement} & \textbf{Rocq name} \\
\midrule
Lemma~\ref{lem:chain}      & chained halving              & \texttt{step\_chain} \\
Prop.~\ref{prop:inv}/\ref{prop:eval} & \textsc{Eval} computes $\Tmap^{q}$ & \texttt{sim\_aux\_correct}, \texttt{sim\_correct} \\
--                         & \texttt{collatz\_sim} spec   & \texttt{collatz\_sim\_correct} \\
Lemma~\ref{lem:affine}     & $2^{q}\Tmap^{q}(n) = A n + C$ & \texttt{affine}, \texttt{A\_odd} \\
Thm.~\ref{thm:shift}       & shift lemma (Terras)         & \texttt{shift}, \texttt{shift\_val} \\
Lemma~\ref{lem:Asmall}     & $A_{q}(r) < 2^{q}$           & \texttt{A\_lt\_pow2} \\
Thm.~\ref{thm:descent}     & class-wide descent           & \texttt{descent} \\
Thm.~\ref{thm:nomin}       & no minimal counterexample    & \texttt{no\_minimal\_counterexample} \\
Prop.~\ref{prop:struct}    & odd / bounded / distinct     & \texttt{Rl\_odd}, \texttt{Rl\_range}, \texttt{Rl\_nodup} \\
Thm.~\ref{thm:cov}         & coverage                     & \texttt{coverage} \\
Thm.~\ref{thm:main}        & sieve soundness              & \texttt{sieve\_sound} \\
Thm.~\ref{thm:main}(c)     & cycles are reported          & \texttt{cycles\_complete} \\
\bottomrule
\end{tabular}
\caption{Correspondence between the paper and \texttt{proofs/CollatzSieve.v}.}
\label{tab:map}
\end{table}

Two conveniences make the development robust. First, \texttt{AC} is unfolded
exactly once, into recursion equations \texttt{Acoef\_S} and
\texttt{Ccoef\_S}; every later proof rewrites with those rather than reducing
\texttt{AC}, so no proof depends on how far \texttt{simpl} chooses to go.
Second, \texttt{sim\_aux} is given the equations \texttt{sim\_aux\_rem0} and
\texttt{sim\_aux\_S} (both closed by \texttt{reflexivity}), because
\texttt{simpl} would otherwise inline the \texttt{let}-bindings and expand
$3n+1$ into its binary representation.

\paragraph{Verification status.}
The file ends with \texttt{Print Assumptions} on every headline result. All
eight report \emph{Closed under the global context}: the development uses no
axioms, no classical logic, no functional extensionality, and contains no
\texttt{Admitted} proofs or \texttt{admit} tactics. It also contains
executable sanity checks closed by \texttt{reflexivity}, for instance
\texttt{Rl 1 = [1; 3]}, \texttt{survives 2 1 = false}, \texttt{Tp 2 1 = 1} and
\texttt{sim 4 3 = 2}. Checked with the Rocq Prover 9.1.1 (OCaml 5.5.0); a
clean build of all 742 lines takes under three seconds.

\paragraph{Reproducing.}
\begin{lstlisting}[style=rocq,numbers=none]
$ rocq compile proofs/CollatzSieve.v
Closed under the global context     (x8)
\end{lstlisting}

% ================================================================= data =====
\section{Computed data}
\label{sec:data}

Table~\ref{tab:data} lists the program's output for $q \leq 24$. The survivor
counts reproduce the classical Terras sieve counts, and the residue sets agree
elementwise with an independent step-by-step implementation of $\Tmap$ that
shares no code with the bit-twiddling version.

\begin{table}[htb]
\centering
\small
\begin{tabular}{@{}rrr@{\hskip 3em}rrr@{}}
\toprule
$q$ & $|\Slev_{q}|$ & eliminated (\%) & $q$ & $|\Slev_{q}|$ & eliminated (\%) \\
\midrule
1 & 1 & 50.0000 & 13 & 367 & 95.5200 \\
2 & 1 & 75.0000 & 14 & 734 & 95.5200 \\
3 & 2 & 75.0000 & 15 & 1295 & 96.0480 \\
4 & 3 & 81.2500 & 16 & 2114 & 96.7743 \\
5 & 4 & 87.5000 & 17 & 4228 & 96.7743 \\
6 & 8 & 87.5000 & 18 & 7495 & 97.1409 \\
7 & 13 & 89.8438 & 19 & 14990 & 97.1409 \\
8 & 19 & 92.5781 & 20 & 27328 & 97.3938 \\
9 & 38 & 92.5781 & 21 & 46611 & 97.7774 \\
10 & 64 & 93.7500 & 22 & 93222 & 97.7774 \\
11 & 128 & 93.7500 & 23 & 168807 & 97.9877 \\
12 & 226 & 94.4824 & 24 & 286581 & 98.2918 \\
\bottomrule
\end{tabular}
\caption{Surviving residue classes and eliminated density by level.}
\label{tab:data}
\end{table}

The only entry ever added to \texttt{cycles} is $1$, contributed at $q=2$
because $\Tmap(1) = 2$ and $\Tmap(2) = 1$. Class $1 \bmod 4$ is therefore
retired at level $2$; by the remark following Theorem~\ref{thm:descent} every
$n \equiv 1 \pmod 4$ with $n > 1$ satisfies $\Tmap^{2}(n) = (3n+1)/4 < n$.

% ================================================================ limits ====
\section{What is \emph{not} proved}
\label{sec:limits}

Theorem~\ref{thm:main} is a soundness result about an elimination procedure.
It is worth being explicit about its reach.

\begin{enumerate}\itemsep4pt
\item \textbf{The sieve cannot prove the conjecture.} Terras~\cite{terras}
      showed the surviving density tends to $0$; Table~\ref{tab:data} exhibits
      that decay. But a set of density zero can be infinite and nonempty, and
      the sieve provides no bound on the surviving classes' members. No
      finite run --- and no limit of runs --- settles Collatz.

\item \textbf{Cycle detection is incidental, not a search.} The equality
      branch fires only when a cycle's minimal element $m$ satisfies
      $m < 2^{L}$ for its $\Tmap$-period $L$ \emph{and} $m$ is still unsieved
      at level $L$. The program is not a cycle-finding algorithm, and
      Theorem~\ref{thm:main}(c) should not be read as one. What the theorem
      does guarantee is the converse direction that matters for soundness: no
      class is discarded on the strength of a fixed point without that fixed
      point being reported.

\item \textbf{Resource growth.} $|\Rlev_{q}| = 2|\Slev_{q-1}|$ grows roughly
      geometrically (empirically by a factor near $1.7$ per level), so memory
      is exponential in $q$. Since the program appends the full survivor list
      to its log at every level, log size is quadratic in the number of levels
      unless suppressed --- which is what the $q \leq 15$ cutoff on line 8 of
      Algorithm~\ref{alg:sieve} is for. No information is lost by that cutoff:
      $|\Slev_{q}|$ is recoverable from the reported percentage $p$ as
      $|\Slev_{q}| = (100 - p)\,2^{q}/100$.

\item \textbf{Machine integers.} The Rocq development reasons about
      mathematical integers. Python's \texttt{int} is arbitrary-precision, so
      this is faithful; a port to fixed-width arithmetic would need an
      overflow argument that is absent here.
\end{enumerate}

% ============================================================ conclusion ====
\section{Conclusion}

The three suspicious-looking features of \texttt{Residual\_Analysis.py} are all
correct, and each for a specific reason. The unconditional $3m+1$ is safe
because oddness is a loop invariant of both loops
(Propositions~\ref{prop:inv} and~\ref{prop:struct}(i)). The truncated shift
computes $\Tmap^{q}(n)$ on the nose, because a truncated division is the tail
of the last $\Tmap$-step rather than a corruption of it
(Proposition~\ref{prop:eval}). And testing the least representative decides
the entire class, because $\Tmap^{q}$ is affine on each class with an odd
multiplier (Theorem~\ref{thm:shift}, Lemma~\ref{lem:Asmall}).

The one defect the analysis surfaced was not in the arithmetic but in the
reporting: the \texttt{cycles} list was populated and never printed. Since
Theorem~\ref{thm:main}(c) identifies that list as precisely the set of
exceptions to the soundness statement, the program's most consequential
possible output was the one it could not emit. That has been corrected.

% ============================================================= references ===
\begin{thebibliography}{9}

\bibitem{terras}
R.~Terras.
\newblock A stopping time problem on the positive integers.
\newblock \emph{Acta Arithmetica}, 30(3):241--252, 1976.

\bibitem{everett}
C.~J.~Everett.
\newblock Iteration of the number-theoretic function $f(2n) = n$,
  $f(2n+1) = 3n+2$.
\newblock \emph{Advances in Mathematics}, 25(1):42--45, 1977.

\bibitem{lagarias}
J.~C.~Lagarias.
\newblock The $3x+1$ problem and its generalizations.
\newblock \emph{The American Mathematical Monthly}, 92(1):3--23, 1985.

\bibitem{lagariasbook}
J.~C.~Lagarias, editor.
\newblock \emph{The Ultimate Challenge: The $3x+1$ Problem}.
\newblock American Mathematical Society, 2010.

\bibitem{tao}
T.~Tao.
\newblock Almost all orbits of the Collatz map attain almost bounded values.
\newblock \emph{Forum of Mathematics, Pi}, 10:e12, 2022.

\bibitem{rocq}
The Rocq Development Team.
\newblock \emph{The Rocq Prover}.
\newblock \url{https://rocq-prover.org/}.

\bibitem{repo}
RobinCodes.
\newblock Collatz --- residual analysis and orbital graphs.
\newblock \url{https://github.com/RobinCodes/Collatz}.

\end{thebibliography}

\end{document}
